% =============================================================================
% DEFINITIVE PLANCK MASS CONVENTION AUDIT FOR DFD
% =============================================================================
\documentclass[11pt]{article}
\usepackage{amsmath,amssymb}
\usepackage[margin=1in]{geometry}
\usepackage{booktabs}
\usepackage{longtable}
\usepackage{tcolorbox}
\usepackage{xcolor}
\usepackage{hyperref}

\newcommand{\MP}{M_P}
\newcommand{\MPbar}{\bar{M}_P}
\newcommand{\GeV}{\,\mathrm{GeV}}
\newcommand{\MeV}{\,\mathrm{MeV}}
\newcommand{\keV}{\,\mathrm{keV}}
\newcommand{\meV}{\,\mathrm{meV}}

\title{Definitive Planck Mass Convention Audit\\for Density Field Dynamics\\[6pt]
\large Resolving the $\MP$ vs.\ $\MPbar$ Confusion in the $\alpha$-Tower}
\author{DFD Convention Audit}
\date{5 April 2026}

\begin{document}
\maketitle

% =============================================================================
\section{The Problem}
% =============================================================================

Two Planck mass conventions exist in the literature:
\begin{align}
  \MP &= \sqrt{\frac{\hbar c}{G}} = 1.220890 \times 10^{19}\GeV
  \qquad\text{(``full'' Planck mass)}, \\[4pt]
  \MPbar &= \frac{\MP}{\sqrt{8\pi}} = 2.435323 \times 10^{18}\GeV
  \qquad\text{(``reduced'' Planck mass)}.
\end{align}
These differ by a factor of $\sqrt{8\pi} \approx 5.013$.

The DFD $\alpha$-tower expresses physical scales as $\MP \alpha^n$ (times
possible prefactors). If $\MP$ is replaced by $\MPbar$ without changing the
formula, every quantity shifts by a factor of $\sim 5$.

This document audits every $\alpha$-tower quantity against the primary source
(DFD v3.3, April 2026) and identifies all inconsistencies in the 20-agent
campaign and associated papers.

% =============================================================================
\section{The Definitive Convention: DFD v3.3}
% =============================================================================

\begin{tcolorbox}[colback=green!5!white, colframe=green!60!black,
  title=\textbf{Result: DFD v3.3 uses the FULL Planck mass $\MP = \sqrt{\hbar c/G}$ throughout}]
Every formula in DFD v3.3 that writes ``$\MP$'' means the \textbf{full}
(unreduced) Planck mass $\MP = 1.2209 \times 10^{19}\GeV$. There is no
instance of the reduced Planck mass in the definitive paper.

The symbol $\MP$ is never redefined to mean $\MPbar$. Every stated numerical
value is consistent with the full Planck mass.
\end{tcolorbox}

\noindent\textbf{Evidence:}
\begin{itemize}
\item $M_R = \MP\alpha^3$: v3.3 states $4.7 \times 10^{12}\GeV$.
  Computation: $1.2209 \times 10^{19} \times (7.297 \times 10^{-3})^3
  = 4.74 \times 10^{12}\GeV$. \checkmark\ (Full $\MP$.)
\item $v = \MP\alpha^8\sqrt{2\pi}$: v3.3 states $246.09\GeV$.
  Computation: $1.2209 \times 10^{19} \times (7.297 \times 10^{-3})^8
  \times 2.507 = 246.09\GeV$. \checkmark\ (Full $\MP$.)
\item $\Lambda_\mathrm{QCD} = \MP\alpha^{19/2}$: v3.3 states $61.20\MeV$.
  Computation: $1.2209 \times 10^{19} \times (7.297 \times 10^{-3})^{9.5}
  = 61.20\MeV$. \checkmark\ (Full $\MP$.)
\item $G\hbar H_0^2/c^5 = \alpha^{57}$: Convention-independent (contains $G$).
  Gives $H_0 = 72.09$ km/s/Mpc regardless of convention choice.
\item $m_3 = \frac{14}{13}\pi\,\MP\alpha^{14}$: v3.3 context implies
  $m_3 \approx 50\meV$. Computation with full $\MP$: $50.2\meV$. \checkmark
\end{itemize}

% =============================================================================
\section{Complete $\alpha$-Tower with the Correct Convention}
% =============================================================================

\begin{tcolorbox}[colback=blue!5!white, colframe=blue!60!black,
  title=\textbf{Convention: $\MP = \sqrt{\hbar c/G} = 1.220890 \times 10^{19}\GeV$ everywhere}]
\end{tcolorbox}

\vspace{0.5cm}

\renewcommand{\arraystretch}{1.5}
\begin{longtable}{@{}p{2.5cm}p{4.5cm}cp{2.5cm}c@{}}
\toprule
\textbf{Quantity} & \textbf{Formula} & \textbf{Exponent} & \textbf{Value} & \textbf{Status} \\
\midrule
\endfirsthead
\toprule
\textbf{Quantity} & \textbf{Formula} & \textbf{Exponent} & \textbf{Value} & \textbf{Status} \\
\midrule
\endhead

$\MP$ (Planck mass)
  & Input
  & $n = 0$
  & $1.2209 \times 10^{19}\GeV$
  & Exact \\
\midrule

$k_\alpha$ (clock coupling)
  & $\alpha^2/(2\pi)$
  & $n = 2$ (dim'less)
  & $8.475 \times 10^{-6}$
  & Theorem \\
\midrule

$M_R$ (Majorana)
  & $\MP\,\alpha^3$
  & $n = 3$
  & $4.74 \times 10^{12}\GeV$
  & Theorem \\
\midrule

$f_\chi$ (ALP scale)
  & $\MP\,\alpha^5$
  & $n = 5$
  & $2.53 \times 10^{8}\GeV$
  & Conjecture \\
\midrule

$v$ (Higgs VEV)
  & $\MP\,\alpha^8\sqrt{2\pi}$
  & $n = 8$
  & $246.09\GeV$
  & Theorem \\
\midrule

$\Lambda_\mathrm{QCD}$
  & $\MP\,\alpha^{19/2}$
  & $n = 19/2$
  & $61.20\MeV$
  & Theorem \\
\midrule

$m_\chi$ (DM mass)
  & $\sqrt{2\pi}\,\MP\,\alpha^{11}$
  & $n = 11$
  & $95.6\keV$
  & Derived \\
\midrule

$m_3$ (neutrino)
  & $\frac{14}{13}\pi\,\MP\,\alpha^{14}$
  & $n = 14$
  & $50.2\meV$
  & Derived \\
\midrule

$H_0$ (Hubble)
  & $G\hbar H_0^2/c^5 = \alpha^{57}$
  & $n = 57$
  & $72.09$ km/s/Mpc
  & Theorem \\
\bottomrule
\end{longtable}

% =============================================================================
\section{Errors Found in the 20-Agent Campaign}
% =============================================================================

\begin{tcolorbox}[colback=red!5!white, colframe=red!70!black,
  title=\textbf{Critical Finding: The R10 and 20-Agent Campaigns Used Mixed Conventions}]
The campaign documents (DEFINITIVE\_SYNTHESIS.md and the Detection Prospects
paper) used the \textit{reduced} Planck mass $\MPbar = 2.435 \times 10^{18}\GeV$
for $f_a$ and $m_\chi$ while writing the symbol ``$\MP$''.
This introduces a factor-of-5 error in some quantities.
\end{tcolorbox}

\subsection{Error 1: $f_a$ value}

The campaign and detection paper state:
\begin{equation}
  f_a \approx 5.04 \times 10^7\GeV.
\end{equation}
But $\MP\alpha^5 = 2.53 \times 10^8\GeV$, not $5.04 \times 10^7\GeV$.

\textbf{What happened:} The value $5.04 \times 10^7$ is $\MPbar\alpha^5$,
i.e.\ the reduced Planck mass was used while the formula was written as
``$\MP\alpha^5$''.

\textbf{Correct value:} $f_\chi = \MP\alpha^5 = 2.53 \times 10^8\GeV$.

\textbf{Impact:} A factor of $\sqrt{8\pi} \approx 5$ in $f_a$.
All photon coupling estimates using $g_{\gamma\chi} = \alpha/(2\pi f_a)$
are affected by a factor of 5. The lifetime estimate
$\tau \propto f_a^2/m_\chi^3$ shifts by a factor of 25.

\subsection{Error 2: $m_\chi$ formula on line 85 of the Synthesis}

The synthesis document (line 85) states:
\begin{equation}
  m_\chi = \sqrt{158} \times \MP \times \alpha^{11} \approx 95.6\keV.
\end{equation}
But $\sqrt{158} \times \MP \times \alpha^{11}
= 12.57 \times 1.2209 \times 10^{19} \times 3.125 \times 10^{-24}
= 479.5\keV$, not $95.6\keV$.

\textbf{What happened:} The formula on line 85 silently uses
$\MPbar$ in place of $\MP$.
With $\MPbar$: $\sqrt{158} \times 2.435 \times 10^{18} \times 3.125 \times 10^{-24}
= 95.65\keV$. \checkmark

\textbf{Two correct ways to write this:}
\begin{align}
  m_\chi &= \sqrt{158}\,\MPbar\,\alpha^{11} \approx 95.65\keV,
  \label{eq:mchi-reduced} \\
  m_\chi &= \sqrt{2\pi}\,\MP\,\alpha^{11} \approx 95.63\keV.
  \label{eq:mchi-full}
\end{align}
These are \textit{not} identical ($\sqrt{158} \neq \sqrt{2\pi}\sqrt{8\pi} = 4\pi$;
$158 \neq 16\pi^2 = 157.914$), but they differ by only 0.027\%, reflecting the
near-coincidence $158 \approx 16\pi^2$.

\textbf{Which is correct?}
The synthesis document itself identifies two distinct origins:
\begin{itemize}
\item \textbf{Line 85:} ``$m_\chi = \sqrt{158} \times \MP \times \alpha^{11}$
  where $158 = 3(k_\mathrm{max} - \dim M_7) - 1$.'' The factor 158 is an
  \textit{exact integer} from $S^3$ spectral geometry.
\item \textbf{Line 216:} ``$m_\chi = \sqrt{2\pi}\,\MP\,\alpha^{11} \approx 96\keV$.''
  This uses the $\sqrt{2\pi}$ normalization inherited from the Higgs VEV.
\end{itemize}

The correct formula with consistent notation is Eq.~\eqref{eq:mchi-full}:
$m_\chi = \sqrt{2\pi}\,\MP\,\alpha^{11}$, using the full Planck mass. This
is consistent with the Higgs formula $v = \MP\alpha^8\sqrt{2\pi}$ and gives
$95.63\keV$.

If Agent 3's spectral geometry result ($V''(0) = 158$ exactly) is to be used,
then the formula must be written as Eq.~\eqref{eq:mchi-reduced}:
$m_\chi = \sqrt{158}\,\MPbar\,\alpha^{11}$, explicitly using $\MPbar$.

\textbf{The near-coincidence} $158 \approx 16\pi^2$ means the two formulas
give numerically indistinguishable results ($95.65$ vs.\ $95.63\keV$).
But they have different physical origins and the notation must be precise.

\subsection{Error 3: ``$\Lambda = \MP\alpha^8 = 19.58\GeV$''}

The R10 campaign sometimes referred to a ``bare'' electroweak scale:
\begin{equation}
  \Lambda = \MP\alpha^8 = 19.58\GeV.
\end{equation}
But $\MP\alpha^8 = 98.17\GeV$, not $19.58\GeV$.
The value $19.58\GeV$ is $\MPbar\alpha^8$.

\textbf{This quantity does not appear in DFD v3.3.} The definitive paper
uses only the Higgs VEV $v = \MP\alpha^8\sqrt{2\pi} = 246.09\GeV$.
The ``bare'' scale $\MP\alpha^8 = 98.17\GeV$ without the $\sqrt{2\pi}$
prefactor has no independent physical role in the published theory.

% =============================================================================
\section{Is There a Single Consistent Convention?}
% =============================================================================

\begin{tcolorbox}[colback=green!5!white, colframe=green!60!black,
  title=\textbf{Answer: Yes. DFD v3.3 uses $\MP = \sqrt{\hbar c/G}$ uniformly.}]
Every formula in the definitive paper uses the \textbf{full} Planck mass.
The reduced Planck mass $\MPbar$ never appears. The $\sqrt{2\pi}$ prefactors
in $v$ and $m_\chi$ are physical normalization factors from the spectral
action, not artifacts of a convention choice.

The convention is:
\begin{equation}
\boxed{\text{All DFD formulas use }\MP = \sqrt{\frac{\hbar c}{G}}
= 1.2209 \times 10^{19}\GeV.}
\end{equation}
\end{tcolorbox}

\subsection{Why the confusion arose}

The reduced Planck mass $\MPbar = \MP/\sqrt{8\pi}$ enters gravitational
physics through the Einstein-Hilbert action:
\begin{equation}
  S_\mathrm{EH} = \frac{\MPbar^2}{2}\int d^4x\,\sqrt{-g}\,R
  = \frac{\MP^2}{16\pi}\int d^4x\,\sqrt{-g}\,R.
\end{equation}
Many particle physics and cosmology references write the Friedmann equation as
$H^2 = \rho/(3\MPbar^2)$, making $\MPbar$ the ``natural'' gravitational scale.

The DFD $\chi$ kinetic term normalization (Section 2.2 of the axion paper)
gives:
\begin{equation}
  f_\chi^2 = \frac{\MP^2}{8\pi}\,\mathrm{Vol}(\mathbb{CP}^2)
  = \MPbar^2\,\mathrm{Vol}(\mathbb{CP}^2).
\end{equation}
This shows that the \textit{natural} normalization of the harmonic 3-form
kinetic term involves $\MPbar$, not $\MP$. The $\alpha^5$ suppression then
acts on $\MPbar$:
\begin{equation}
  f_\chi = \MPbar\,\alpha^5 \times \sqrt{\mathrm{Vol}(\mathbb{CP}^2)}
  \approx 5.04 \times 10^7\GeV \times O(1).
\end{equation}

This is the source of the confusion: the kinetic normalization
$\MP^2/(8\pi) = \MPbar^2$ naturally produces $\MPbar$, and the campaign
agents used this value while still writing ``$\MP$''.

\subsection{Resolution: two equivalent notations}

For any quantity, one can write:
\begin{equation}
  X = \MP\,\alpha^n \times (\text{prefactor with }\MP)
  = \MPbar\,\alpha^n \times (\text{prefactor with }\MPbar}).
\end{equation}
The DFD convention is to always use $\MP$ and absorb any $1/\sqrt{8\pi}$
factors into the prefactor. For example:
\begin{equation}
  m_\chi = \sqrt{2\pi}\,\MP\,\alpha^{11}
  = \sqrt{2\pi} \times \sqrt{8\pi} \times \MPbar\,\alpha^{11}
  = 4\pi\,\MPbar\,\alpha^{11}.
\end{equation}
Both are correct; the DFD convention uses the first form.

% =============================================================================
\section{Corrected Quantities for the Campaign}
% =============================================================================

The following campaign quantities must be corrected:

\renewcommand{\arraystretch}{1.4}
\begin{longtable}{@{}p{2.5cm}p{3.5cm}p{3.5cm}p{3.2cm}@{}}
\toprule
\textbf{Quantity} & \textbf{Campaign (wrong)} & \textbf{Corrected} & \textbf{Shift} \\
\midrule

$f_a$
  & $5.04 \times 10^7\GeV$
  & $2.53 \times 10^8\GeV$
  & $\times 5.01$ \\
\midrule

$g_{\gamma\chi}$
  & $2.3 \times 10^{-11}\,\mathrm{GeV}^{-1}$
  & $4.6 \times 10^{-12}\,\mathrm{GeV}^{-1}$
  & $\div 5.01$ \\
\midrule

$\tau(\chi\to\gamma\gamma)$
  & $\sim 300$ s
  & $\sim 7500$ s
  & $\times 25$
  \\
  & \multicolumn{3}{l}{\footnotesize (Still catastrophically short; does not resolve the lifetime crisis)} \\
\midrule

$\Omega_\chi h^2$ (misal.)
  & $6.33 \times 10^{-7}\,\theta_i^2$
  & $1.59 \times 10^{-5}\,\theta_i^2$
  & $\times 25$
  \\
  & \multicolumn{3}{l}{\footnotesize (Still 4 orders of magnitude too small; misalignment still fails)} \\
\midrule

``$\MP\alpha^8$''
  & $19.58\GeV$
  & $98.17\GeV$
  & $\times 5.01$ \\
\bottomrule
\end{longtable}

\textbf{Quantities that are UNAFFECTED:}
\begin{itemize}
\item $m_\chi \approx 95.6\keV$: The numerical value is correct in all documents
  (the convention error in the formula cancels between the mass scale and the
  prefactor).
\item $v = 246.09\GeV$: Always stated correctly.
\item $\Lambda_\mathrm{QCD} = 61.20\MeV$: Always stated correctly.
\item $M_R = 4.74 \times 10^{12}\GeV$: Always stated correctly.
\item $H_0 = 72.09$ km/s/Mpc: Convention-independent formula.
\item $\Omega_\chi/\Omega_b = 16/3$: Independent of $f_a$.
\end{itemize}

% =============================================================================
\section{The Prefactor Question: $\sqrt{158}$ vs.\ $\sqrt{2\pi}$}
% =============================================================================

The synthesis document contains two formulas for $m_\chi$ that are
\textit{inconsistent in their prefactors}:

\begin{align}
  \text{Line 85:}\quad m_\chi &= \sqrt{158}\,\MP\,\alpha^{11}
  \quad\text{(should be }\MPbar\text{)}, \\
  \text{Line 216:}\quad m_\chi &= \sqrt{2\pi}\,\MP\,\alpha^{11}
  \quad\text{(correct with full }\MP\text{)}.
\end{align}

These arise from different physical derivations:
\begin{itemize}
\item $\sqrt{158}$: From Agent 3's spectral geometry result.
  $V''(0) = 158 = 3(k_\mathrm{max} - \dim M_7) - 1$ is an \textit{exact integer}
  from the $S^3$ heat kernel. With the reduced mass,
  $m_\chi = \sqrt{158}\,\MPbar\,\alpha^{11} = 95.65\keV$.

\item $\sqrt{2\pi}$: From the Higgs-sector normalization inherited through the
  see-saw structure. With the full mass,
  $m_\chi = \sqrt{2\pi}\,\MP\,\alpha^{11} = 95.63\keV$.
\end{itemize}

The near-coincidence arises because:
\begin{equation}
  158 \approx 16\pi^2 = (\sqrt{2\pi} \times \sqrt{8\pi})^2
  = (2\pi)({8\pi}) \times \frac{158}{16\pi^2}
  \approx 157.914 \times 1.00055.
\end{equation}

\textbf{These formulas must be distinguished:}
\begin{itemize}
\item If $158$ is exact (spectral geometry): use $\sqrt{158}\,\MPbar\,\alpha^{11}$.
\item If the $\sqrt{2\pi}$ normalization is fundamental: use
  $\sqrt{2\pi}\,\MP\,\alpha^{11}$.
\item They differ by $0.027\%$ --- well below any foreseeable experimental
  sensitivity, but the \textit{theoretical} distinction matters for the
  derivation chain.
\end{itemize}

% =============================================================================
\section{Recommended Convention Lock}
% =============================================================================

\begin{tcolorbox}[colback=blue!5!white, colframe=blue!60!black,
  title=\textbf{DFD Convention Lock (to be applied in all future documents)}]

\begin{enumerate}
\item \textbf{Symbol $\MP$} always means $\sqrt{\hbar c/G} = 1.2209 \times 10^{19}\GeV$.
\item \textbf{Symbol $\MPbar$} always means $\MP/\sqrt{8\pi} = 2.4353 \times 10^{18}\GeV$.
\item If $\MPbar$ is used, it must be \textit{written explicitly} as $\MPbar$.
  It must never be written as $\MP$.
\item \textbf{The canonical $\alpha$-tower uses $\MP$:}
  \begin{align}
    M_R &= \MP\,\alpha^3, \\
    f_\chi &= \MP\,\alpha^5, \\
    v &= \MP\,\alpha^8\,\sqrt{2\pi}, \\
    \Lambda_\mathrm{QCD} &= \MP\,\alpha^{19/2}, \\
    m_\chi &= \sqrt{2\pi}\,\MP\,\alpha^{11}, \\
    m_3 &= \tfrac{14}{13}\pi\,\MP\,\alpha^{14}, \\
    G\hbar H_0^2/c^5 &= \alpha^{57}.
  \end{align}
\item \textbf{Numerical cross-check}: any stated value must be verified
  against $\MP = 1.2209 \times 10^{19}\GeV$, not $2.4353 \times 10^{18}\GeV$.
\item \textbf{When the reduced mass arises naturally} (e.g.\ from kinetic
  normalization $\MP^2/(8\pi)$), write the formula in terms of $\MP$ with
  explicit $1/\sqrt{8\pi}$ factors, or use $\MPbar$ with the bar notation.
\end{enumerate}
\end{tcolorbox}

% =============================================================================
\section{Impact on Published Results}
% =============================================================================

\subsection{What changes}

\begin{itemize}
\item \textbf{$f_\chi = 2.53 \times 10^8\GeV$, not $5.04 \times 10^7\GeV$.}
  This is 5 times larger than the campaign value. The ALP is less strongly
  coupled to the SM by a factor of 5.

\item \textbf{$g_{\gamma\chi} = 4.6 \times 10^{-12}\,\mathrm{GeV}^{-1}$,
  not $2.3 \times 10^{-11}\,\mathrm{GeV}^{-1}$.}
  The photon coupling drops by a factor of 5.

\item \textbf{$\tau(\chi\to\gamma\gamma) \sim 7500$ s, not $\sim 300$ s.}
  Still catastrophically short (2 hours vs.\ 5 minutes). The lifetime crisis
  is not resolved by the convention correction. The Z$_2$ stabilization
  question remains critical.

\item \textbf{Misalignment relic density increases by $\times 25$.}
  Still $\sim 10^4$ too small. The production mechanism problem is not
  resolved by the convention correction.
\end{itemize}

\subsection{What does NOT change}

\begin{itemize}
\item $m_\chi \approx 95.6\keV$ (value correct in all documents)
\item $v = 246.09\GeV$ (always correct)
\item $\Lambda_\mathrm{QCD} = 61.20\MeV$ (always correct)
\item $H_0 = 72.09$ km/s/Mpc (convention-independent)
\item $\Omega_\chi/\Omega_b = 16/3$ (independent of $f_a$)
\item All CMB-S4 and ELT/ANDES predictions (unaffected)
\end{itemize}

% =============================================================================
\section{Numerical Reference Table}
% =============================================================================

For quick reference, all $\alpha$-tower quantities with both conventions:

\renewcommand{\arraystretch}{1.3}
\begin{longtable}{@{}lccc@{}}
\toprule
\textbf{Formula} & \textbf{With $\MP$} & \textbf{With $\MPbar$}
  & \textbf{DFD uses} \\
\midrule
$\alpha^3 \times \text{mass}$ & $4.74 \times 10^{12}\GeV$ & $9.46 \times 10^{11}\GeV$ & $\MP$ \\
$\alpha^5 \times \text{mass}$ & $2.53 \times 10^{8}\GeV$ & $5.04 \times 10^{7}\GeV$ & $\MP$ \\
$\alpha^8 \times \text{mass}$ & $98.17\GeV$ & $19.58\GeV$ & $\MP$ \\
$\alpha^8 \sqrt{2\pi} \times \text{mass}$ & $246.09\GeV$ & $49.09\GeV$ & $\MP$ \\
$\alpha^{19/2} \times \text{mass}$ & $61.20\MeV$ & $12.21\MeV$ & $\MP$ \\
$\sqrt{2\pi}\,\alpha^{11} \times \text{mass}$ & $95.63\keV$ & $19.07\keV$ & $\MP$ \\
$\sqrt{158}\,\alpha^{11} \times \text{mass}$ & $479.5\keV$ & $95.65\keV$ & $\MPbar$ \\
$\frac{14}{13}\pi\,\alpha^{14} \times \text{mass}$ & $50.2\meV$ & $10.0\meV$ & $\MP$ \\
\bottomrule
\end{longtable}

\end{document}
